\section{Kokkuvõte} \label{kokkuvote}

Käesolev bakalaureusetöö uuris tehisintellekti rakendamise teostatavust eestikeelse tehnilise informaatika lõputöö genereerimisel. Uurimuse käigus võrreldi nelja juhtivat suurt keelemudelit koos nende agentsete tööriistadega -- Google Gemini 3.1 Pro (Antigravity), OpenAI GPT-5.4 (ChatGPT vestlusliides ja ChatGPT Codex laiendus Visual Studio Code'is), Anthropic Claude Sonnet 4.6 (Cowork) ja Anthropic Claude Opus 4.6 (Cowork) -- nende võimekuse osas luua struktureeritud, akadeemiliselt korrektset ja sisuliselt väärtuslikku eestikeelset teksti. Iga kombinatsioon sai ühe lühikese eestikeelse kasutajaviibega ülesande genereerida terviklik LaTeX-vormingus bakalaureusetöö Tartu Ülikooli arvutiteaduse instituudi nõuetele vastavalt, valides ise informaatikaga seotud teema ning kasutades töökausta paigutatud UniTartuCS malli. Cowork on sama ettevõtte Claude Code'ist eraldiseisev toode. Iga kombinatsiooni kohta tehti üks jooks, mistõttu tulemusi tuleb käsitleda eksploratiivselt -- need iseloomustavad konkreetset katset, mitte üldist mudelivõimekust.

Selle katse põhjal saab järeldada, et TI mudelid koos agentsete tööriistadega suudavad genereerida rahuldava kuni hea kvaliteediga eestikeelset akadeemilist teksti isegi ühe kasutajaviibega, ilma iteratiivse juhendamise või täiendava konteksti lisamiseta. Eriti silmapaistev oli see, et Claude Opus 4.6 + Cowork kombinatsioon genereeris 61-leheküljelise tehniliselt sügava eestikeelse lõputöö ilma ühegi täiendava juhiseta.

Nelja kombinatsiooni võrdluses joonistusid välja fundamentaalselt erinevad käitumismustrid. Claude Opus 4.6 + Cowork genereeris kõige mahukama väljundi -- 61-leheküljelise eestikeelse lõputöö paroolihalduri teemal, mis oli tehniliselt sügav ja strukturaalselt terviklik, sisaldas 35 viidet (tsitaatide tihedus 0,74 lehekülje kohta) ning vaid 5 kirjaviga terve töö peale. Claude Sonnet 4.6 + Cowork genereeris 40-leheküljelise eestikeelse lõputöö TI-põhise lõputöö genereerimise teemal, paistes silma eriti puhta eestikeelse stiiliga (vaid 0,13 anglitsismi lehekülje kohta). GPT-5.4 keeldus tavapärase ChatGPT vestlusliidese kaudu terviklikku tööd kirjutamast, viidates akadeemilise aususe piirangutele, ja pakkus selle asemel 8-leheküljelise ingliskeelse teostatavusjuhendi. ChatGPT Codex laienduse kaudu Visual Studio Code'is genereeris sama mudel sama viipega 47-leheküljelise eestikeelse lõputöö RAG-põhise küsimus-vastus süsteemi teemal, seejuures vaid ühe tuvastatud kirjaveaga. Gemini 3.1 Pro + Antigravity kombinatsioon genereeris 39-leheküljelise eestikeelse töö generatiivse TI mõjust programmeerimisele, kuid sisaldas kõige rohkem kirjavigu (28 viga ehk 0,72 viga lehekülje kohta), väikseimat viidete arvu (9 allikat) ning -- kõige olulisemana -- täielikult fabritseeritud empiirilise osa (30 väljamõeldud tudengit, fiktiivsed ettevõtted ja statistilised tulemused), mida käsitleme akadeemilise väärkasutuse kõrge riskiga juhtumina. Need tulemused näitavad, et mudel-tööriista kombinatsioonide juhiste järgimise võimekus, keelevalik, teemale vastavus ja ortograafiline täpsus erinevad drastiliselt.

Eksperiment kinnitas ka agentse liidese kriitilist rolli: Anthropicu Cowork, Google'i Antigravity ja ChatGPT Codex andsid oluliselt erineva tulemuse kui puhas vestlusliides -- kõige ilmekamalt GPT-5.4 puhul, kus sama mudel ja sama viibega anti vestlusliideses keeldumine ja Codex-laiendusega terviklik eestikeelne dokument. See erinevus ei ole tõlgendatav kui ohutus\-piirangu ``ületamine'', vaid näitab, et ``mudelikäitumine'' on tegelikult mudeli, liidese, süsteemijuhiste ja ohutuspoliitika koosmõju.

Kõige tõsisem probleem oli viidete hallutsineerimine -- kombinatsioonid leiutasid olematud allikad 15--25\% juhtudest (Opus 4.6 + Cowork ca 15\%, Sonnet 4.6 + Cowork ca 20\%, GPT-5.4 + Codex ca 20\%, Gemini 3.1 Pro + Antigravity ca 25\%). See tähendab, et iga TI-genereeritud viide vajab inimese poolt kontrollimist. Samuti jäi TI-genereeritud teksti sisuline sügavus ja originaalsus oodatust madalamaks ning Gemini 3.1 Pro + Antigravity puhul oli kogu empiiriline uurimus fabritseeritud -- tudengeid, andmeid ja eksperimenti reaalselt ei eksisteerinud. Ka Opus 4.6 + Cowork väljund sisaldas empiirilisi väiteid (149 testi, SUS skoor 77,0 jms), mida käesoleva uurimuse autor iseseisvalt ei verifitseerinud ja mis tuleb samuti lugeda kontrollimata väideteks.

TI kasutamise majanduslik kulu on Sonnet-klassi kombinatsioonidega mõistlik: \$23--42 API kaudu või \$20--40 kuutellimuste kaudu. Opus 4.6 + Cowork on oluliselt kallim (\$160--220 API kaudu, Claude Max kuutellimus \$100--200/kuu), kuid andis selles katses kõrgeima kvaliteedi. Mõlemad variandid on soodsamad kui professionaalse toimetamise teenus.

Uurimuse põhjal saab järeldada, et TI mudeleid saab efektiivselt kasutada lõputöö kirjutamise abivahendina, kuid mitte asendajana. TI tugevused -- struktureerimine, mustandi loomine, keeleline toimetamine ja LaTeX-vormindamine -- säästavad üliõpilasele märkimisväärset aega. Samas on inimese panus endiselt asendamatu sisulise analüüsi, originaalsete järelduste, viidete kontrollimise ja valdkonnaspetsiifilise terminoloogia osas.

Selle katse põhjal soovitame kasutada Claude Opus 4.6 mudelit koos Anthropicu Cowork'i agentse tööriistaga, eriti terviklike peatükkide ja keeruliste akadeemiliste tekstide genereerimisel; sarnase ülesehitusega Claude Code sobib samaks otstarbeks. Piiratud eelarve korral oli Claude Sonnet 4.6 + Cowork selle katse parim hinna-kvaliteedi suhtega valik. GPT-5.4 sobib hästi keelelise korrektuuri jaoks ja lühikeste tekstilõikude genereerimiseks; tervikliku töö genereerimisel tuleb arvestada, et ChatGPT vestlusliides võib ülesandest keelduda, samas kui ChatGPT Codex laiendus tootis sama viipega tervikliku eestikeelse dokumendi (vt peatükk~\ref{subsec:codex}). Ideaalne on agentse tööriista kasutamine koos Opus või Sonnet mudeliga (vt peatükk~\ref{claudecode}). Kuna iga kombinatsiooni kohta tehti üks jooks, tuleks neid soovitusi võtta eksploratiivse viitena, mitte lõpliku pingereana.

Käesolev töö panustab eestikeelsesse TI-uurimusse, pakkudes esimese süstemaatilise ühe kasutajaviibega (\textit{single-prompt}) võrdluse kaasaegsete mudel-tööriista kombinatsioonide suutlikkusest tervikliku eestikeelse bakalaureusetöö genereerimisel. Tulemused on kasulikud nii üliõpilastele, kes soovivad TI-d vastutustundlikumalt kasutada, kui ka õppejõududele ja ülikoolidele, kes kujundavad TI kasutamise poliitikat akadeemilises kontekstis.

Edaspidine uurimistöö võiks keskenduda mitmele suunale. Üks võimalik fookus on spetsialiseeritud akadeemiliste TI tööriistade hindamine ning eesti keelele kohandatud mudelite (nt EstLLM projekti \cite{dorkin2026estllm} baasil loodud variantide) võrdlemine suurte universaalsete mudelitega. Teine oluline uurimissuund on RAG-süsteemide mõju hindamine viidete kvaliteedile ning iteratiivse mitmeetapilise promptimise ja ühe viibega genereerimise kvaliteedi süstemaatiline võrdlus. Samuti on oluline samu kombinatsioone korduvalt testida, et hinnata mitteterministlikkusest tulenevat varieeruvust. Pikas perspektiivis tasub hinnata ka TI kasutamise mõju üliõpilaste akadeemilisele arengule.
